\subsection{Incremental population result}\label{sec:incremental-results}

The separate \FrontierIncrementSources{}-source study compares the extension
with the base controller. Sources follow the same independent uniform-deal
and seat law as Section~\ref{sec:uniform-results}, with 1,024 matches in each
parent--opponent population. Every source remains in the denominator. The
comparison uses the exact increment from independent terminal replay at
adopted histories, and zero where the complete policies coincide.
The fixed primary criterion requires positive simultaneous lower bounds for
both equal-opponent parent mixtures. Table~\ref{tab:frontier-increment} reports
the complete cohort; Appendix~\ref{app:frontier} records the scope correction,
all source outcomes and computational work.

\begin{table}[htbp]
\centering\small
% Complete incremental study: hybrid minus existing controller.
\begin{tabular}{llrrr}
\toprule
Parent & Opponent & Sources & Mean (pp) & Lower bound (pp) \\
\midrule
1 & WP & 1,024 & $6.77034\times 10^{-4}$ & $6.64356\times 10^{-5}$ \\
1 & Point & 1,024 & 0.00189203 & $2.73724\times 10^{-5}$ \\
2 & WP & 1,024 & 0.00106778 & $1.10018\times 10^{-5}$ \\
2 & Point & 1,024 & $5.00209\times 10^{-4}$ & $7.94267\times 10^{-6}$ \\
\midrule
1 & Equal mixture & 2,048 & 0.00128453 & $4.69040\times 10^{-5}$ \\
2 & Equal mixture & 2,048 & $7.83993\times 10^{-4}$ & $9.47225\times 10^{-6}$ \\
\bottomrule
\end{tabular}

\caption{Frontier extension minus base controller: expected terminal win-score
increments in percentage points. Bounds are simultaneous at 95\% across the
four populations; the fixed equal-weight mixtures inherit that coverage.
All 4,096 sources contribute, including missing roots and unchanged policies.
Lower bounds are rounded downward.}
\label{tab:frontier-increment}
\end{table}

The primary criterion is \FrontierIncrementCriterion{}. All four individual
population lower bounds are also positive. The extension improves 26 source
beliefs, with equal-opponent means of \FrontierIncrementMeanOne{} and
\FrontierIncrementMeanTwo{} percentage points. Thus a failed root reference
can conceal a recoverable improvement in naturally reached play. The measured
increment is small, and realized match-score differences are zero in every
population. The earlier base-controller 0.1-point criterion remains a separate
failed test. Additional planning uses 111,874,298 physical transitions across
this cohort, within the unchanged maximum per-root allowance.
